\documentclass[11pt,a4paper]{article}
%%%%%%%%%%%%%%%%%%%%%%%%% Credit %%%%%%%%%%%%%%%%%%%%%%%%

% template ini dibuat oleh martin.manullang@if.itera.ac.id untuk dipergunakan oleh seluruh sivitas akademik itera.


%%%%%%%%%%%%%%%%%%%%%%%%% Credit %%%%%%%%%%%%%%%%%%%%%%%%

% template ini dibuat oleh martin.manullang@if.itera.ac.id untuk dipergunakan oleh seluruh sivitas akademik itera.

%%%%%%%%%%%%%%%%%%%%%%%%% PACKAGE starts HERE %%%%%%%%%%%%%%%%%%%%%%%%
\usepackage{graphicx}
\usepackage{caption}
\usepackage[expansion=false]{microtype}
\captionsetup[table]{name=Tabel}
\captionsetup[figure]{name=Gambar}
\usepackage{tabulary}
\usepackage{minted}
% \usepackage{amsmath}
\usepackage{fancyhdr}
% \usepackage{amssymb}
% \usepackage{amsthm}
\usepackage{placeins}
% \usepackage{amsfonts}
\usepackage{graphicx}
\usepackage[all]{xy}
\usepackage{tikz}
\usepackage{verbatim}
\usepackage[left=2cm,right=2cm,top=3cm,bottom=2.5cm]{geometry}
\usepackage{hyperref}
\hypersetup{
    colorlinks,
    linkcolor={red!50!black},
    citecolor={blue!50!black},
    urlcolor={blue!80!black}
}
\usepackage{caption}
\usepackage{subcaption}
\usepackage{multirow}
\usepackage{psfrag}
\usepackage[T1]{fontenc}
\usepackage[scaled]{beramono}
% Enable inserting code into the document
\usepackage{listings}
\usepackage{xcolor} 
% custom color & style for listing
\definecolor{codegreen}{rgb}{0,0.6,0}
\definecolor{codegray}{rgb}{0.5,0.5,0.5}
\definecolor{codepurple}{rgb}{0.58,0,0.82}
\definecolor{backcolour}{rgb}{0.95,0.95,0.92}
\definecolor{LightGray}{gray}{0.9}
\lstdefinestyle{mystyle}{
	backgroundcolor=\color{backcolour},   
	commentstyle=\color{green},
	keywordstyle=\color{codegreen},
	numberstyle=\tiny\color{codegray},
	stringstyle=\color{codepurple},
	basicstyle=\ttfamily\footnotesize,
	breakatwhitespace=false,         
	breaklines=true,                 
	captionpos=b,                    
	keepspaces=true,                 
	numbers=left,                    
	numbersep=5pt,                  
	showspaces=false,                
	showstringspaces=false,
	showtabs=false,                  
	tabsize=2
}
\lstset{style=mystyle}
\renewcommand{\lstlistingname}{Kode}
%%%%%%%%%%%%%%%%%%%%%%%%% PACKAGE ends HERE %%%%%%%%%%%%%%%%%%%%%%%%

%%%%%%%%%%%%%%%%%%%%%%%%% Data Diri %%%%%%%%%%%%%%%%%%%%%%%%
\newcommand{\student}{\textbf{Joshia Fernandes Sectio Purba (122140170)}}
\newcommand{\course}{\textbf{Deep Learning}}
\newcommand{\assignment}{\textbf{
Tugas Eksplorasi Vision Transformer}} % Menambahkan kurung tutup yang hilang

%%%%%%%%%%%%%%%%%%% using theorem style %%%%%%%%%%%%%%%%%%%%
\newtheorem{thm}{Theorem}
\newtheorem{lem}[thm]{Lemma}
\newtheorem{defn}[thm]{Definition}
\newtheorem{exa}[thm]{Example}
\newtheorem{rem}[thm]{Remark}
\newtheorem{coro}[thm]{Corollary}
\newtheorem{quest}{Question}[section]
%%%%%%%%%%%%%%%%%%%%%%%%%%%%%%%%%%%%%%%%
\usepackage{lipsum}%% a garbage package you don't need except to create examples.
\usepackage{fancyhdr}
\pagestyle{fancy}
\lhead{Joshia Fernandes Sectio Purba (122140170}
\rhead{ \thepage}
\cfoot{\textbf{J
Tugas Eksplorasi Vision Transformer}}
\renewcommand{\headrulewidth}{0.4pt}
\renewcommand{\footrulewidth}{0.4pt}

%%%%%%%%%%%%%%  Shortcut for usual set of numbers  %%%%%%%%%%%

\newcommand{\N}{\mathbb{N}}
\newcommand{\Z}{\mathbb{Z}}
\newcommand{\Q}{\mathbb{Q}}
\newcommand{\R}{\mathbb{R}}
\newcommand{\C}{\mathbb{C}}
\setlength\headheight{14pt}

%%%%%%%%%%%%%%%%%%%%%%%%%%%%%%%%%%%%%%%%%%%%%%%%%%%%%%%555
\begin{document}
\thispagestyle{empty}
\begin{center}
	\includegraphics[scale = 0.15]{Figure/ifitera-header.png}
	\vspace{0.1cm}
\end{center}
\noindent
\rule{17cm}{0.2cm}\\[0.3cm]
Nama: \student \hfill Tugas Ke: \assignment\\[0.1cm]
Mata Kuliah: \course \hfill Tanggal: 22 November 2025\\
\rule{17cm}{0.05cm}
\vspace{0.1cm}



%%%%%%%%%%%%%%%%%%%%%%%%%%%%%%%%%%%%%%%%%%%%% BODY DOCUMENT %%%%%%%%%%%%%%%%%%%%%%%%%%%%%%%%%%%%%%%%%%%%%
\textbf{Repository:} \url{hhttps://github.com/Joshia05/Tugas-deep-learning1.git}

\section{PENDAHULUAN}

\subsection{Latar Belakang}
Pengenalan citra merupakan salah satu tantangan utama dalam bidang visi komputer, dengan aplikasi yang luas, termasuk dalam pengenalan makanan. Model Vision Transformer (ViT) telah menunjukkan hasil yang sangat baik dalam klasifikasi citra dengan memanfaatkan mekanisme self-attention untuk menangkap hubungan global antar bagian gambar. Namun, ViT membutuhkan dataset besar dan sumber daya komputasi yang tinggi, yang membatasi penggunaannya dalam beberapa aplikasi.

Sebagai alternatif, Data-efficient Transformer (DeiT) dikembangkan untuk meningkatkan efisiensi pelatihan dan inferensi tanpa mengorbankan akurasi. Dengan memanfaatkan teknik distillation, DeiT mampu mencapai performa yang kompetitif meskipun menggunakan dataset yang lebih kecil dan lebih sedikit sumber daya komputasi. Percobaan ini bertujuan untuk membandingkan kedua model dalam klasifikasi citra makanan, dengan fokus pada akurasi, efisiensi komputasi, dan perbandingan performa antara ViT dan DeiT, untuk menentukan model yang lebih sesuai digunakan dalam aplikasi praktis.

\subsection{Motivasi Perbandingan Model}
Motivasi perbandingan antara Vision Transformer (ViT) dan Data-efficient Transformer (DeiT) dalam percobaan ini didorong oleh kebutuhan untuk mengeksplorasi model yang dapat memberikan akurasi tinggi dalam pengenalan citra, khususnya dalam klasifikasi makanan, sambil mempertimbangkan efisiensi komputasi. Meskipun kedua model ini menggunakan arsitektur berbasis transformer, ViT dikenal membutuhkan dataset besar dan komputasi tinggi untuk dapat mencapai performa optimal. Oleh karena itu, ViT cenderung lebih mahal dalam hal penggunaan sumber daya, yang membatasi penggunaannya dalam aplikasi dengan keterbatasan komputasi. Di sisi lain, DeiT menawarkan efisiensi yang lebih baik dengan teknik distillation, memungkinkan model lebih ringan dan cepat, namun tetap mempertahankan akurasi yang sangat baik.

Perbandingan ini penting untuk menentukan model yang paling sesuai dengan berbagai kondisi dan kebutuhan aplikasi pengenalan citra. Misalnya, untuk aplikasi yang membutuhkan sumber daya terbatas, seperti perangkat mobile atau edge computing, DeiT adalah pilihan yang lebih efisien. Sementara itu, jika akurasi maksimal diperlukan dan sumber daya yang cukup tersedia, ViT dapat menjadi pilihan yang lebih baik. Percobaan ini bertujuan untuk mengevaluasi kedua model berdasarkan akurasinya serta kecepatan inferensinya, sehingga hasilnya dapat digunakan untuk memilih model yang paling sesuai dalam konteks pengenalan citra makanan dengan mempertimbangkan faktor-faktor tersebut.

\subsection{Tujuan Eksperimen}
Tujuan eksperimen ini adalah untuk membandingkan dua model, ViT dan DeiT, dalam hal:
\begin{itemize}
    \item Akurasi klasifikasi citra makanan.
    \item Efisiensi komputasi, termasuk waktu inferensi dan penggunaan memori.
    \item Visualisasi perbandingan performa dalam berbagai aspek.
\end{itemize}

\section{LANDASAN TEORI}

\section{LANDASAN TEORI}

\subsection{Transformer dan Self-Attention}
Transformer merupakan model berbasis perhatian yang pertama kali diperkenalkan dalam bidang pemrosesan bahasa alami (Natural Language Processing/NLP), namun seiring dengan perkembangan teknologi, model ini juga telah diterapkan dengan sangat baik dalam bidang pengolahan citra. Transformer menggunakan mekanisme **self-attention** untuk menangkap ketergantungan jangka panjang antar bagian dalam data. Dalam pengolahan citra, **self-attention** memungkinkan model untuk memahami hubungan global antar bagian citra, yang membuatnya efektif dalam tugas-tugas seperti klasifikasi citra dan deteksi objek.

Namun, meskipun sangat efektif dalam menangkap informasi kontekstual dalam citra, model Transformer seperti **Vision Transformer (ViT)** cenderung memerlukan **dataset besar** dan **komputasi tinggi**. Kebutuhan akan sumber daya yang besar inilah yang menjadi tantangan utama dalam penerapan ViT pada aplikasi dengan keterbatasan sumber daya komputasi, seperti perangkat mobile dan aplikasi berbasis edge computing.

\subsection{Vision Transformer (ViT)}
**Vision Transformer (ViT)** adalah salah satu model yang mengadaptasi arsitektur Transformer untuk tugas klasifikasi citra. ViT membagi gambar menjadi beberapa **patches** dan menerapkan mekanisme **self-attention** pada patch tersebut untuk menangkap ketergantungan antar bagian gambar. Kelebihan utama ViT terletak pada kemampuannya untuk menangkap hubungan global dalam citra, yang memungkinkan model ini mengklasifikasikan citra dengan akurasi yang tinggi. Namun, kekurangan utama ViT adalah ketergantungannya pada **dataset besar** dan **komputasi yang sangat tinggi**, yang membatasi penggunaannya pada aplikasi dengan sumber daya terbatas.

\subsection{Data-efficient Transformer (DeiT)}
**Data-efficient Transformer (DeiT)** adalah variasi dari ViT yang dirancang untuk meningkatkan efisiensi pelatihan dengan menggunakan teknik **distillation**. Dalam DeiT, model kecil (student) dilatih dengan memanfaatkan model besar (teacher), yang membantu mengurangi kebutuhan data besar dan komputasi tinggi. Meskipun akurasinya sedikit lebih rendah dibandingkan ViT, DeiT menawarkan **efisiensi komputasi** yang lebih baik, menjadikannya pilihan yang lebih sesuai untuk aplikasi yang membutuhkan **kecepatan inferensi** dan **penggunaan sumber daya yang lebih efisien**. Keunggulan DeiT adalah kemampuannya untuk **memproses data lebih cepat** dengan akurasi yang tetap kompetitif, meskipun menggunakan dataset yang lebih kecil dan lebih sedikit sumber daya.

\subsection{Perbandingan Model ViT dan DeiT}
Meskipun ViT dan DeiT keduanya menggunakan mekanisme **self-attention** yang sama, perbedaan mendasar antara keduanya terletak pada cara mereka menangani **data** dan **komputasi**. ViT memanfaatkan **attention global** yang memerlukan banyak data dan komputasi untuk berfungsi dengan baik, sedangkan DeiT menggunakan teknik distillation untuk mengurangi kebutuhan akan data besar, menjadikannya lebih efisien dalam hal waktu pelatihan dan inferensi. Dalam percobaan ini, kedua model ini dibandingkan untuk menilai **akurasi klasifikasi citra**, **efisiensi komputasi**, dan **waktu inferensi**, khususnya dalam konteks **klasifikasi citra makanan**. Dengan menguji kedua model ini, kami berharap dapat memilih model yang lebih cocok untuk aplikasi dengan **sumber daya terbatas** namun tetap memprioritaskan **akuratnya hasil klasifikasi**.


\subsection{Perbedaan Kunci}
\begin{table}[h]
\caption{Perbandingan Teoritis Model ViT vs DeiT}
\label{tab:theory-comparison}
\centering
\begin{tabular}{|l|l|l|}
\hline
\textbf{Aspek} & \textbf{ViT} & \textbf{DeiT} \\ \hline
Attention Mechanism & Global Attention & Distilled Global Attention \\ \hline
Training Data Requirement & Large Dataset & Small Dataset (with distillation) \\ \hline
Performance Efficiency & Lower Efficiency & More Efficient with distillation \\ \hline
\end{tabular}
\end{table}

\section{METODOLOGI}

\subsection{Deskripsi Dataset}
Dataset yang digunakan dalam eksperimen ini terdiri dari **gambar makanan**, yang terbagi menjadi dua bagian: **838 gambar untuk pelatihan** dan **270 gambar untuk pengujian**. Setiap gambar pada dataset ini diklasifikasikan ke dalam **5 kategori makanan** yang berbeda, yang mencakup beragam jenis makanan. Dataset ini dirancang untuk menguji kemampuan model dalam melakukan **klasifikasi citra** dengan berbagai variasi dalam kategori makanan yang berbeda. Proses pengumpulan data bertujuan untuk memastikan keragaman gambar yang mencakup berbagai kondisi pencahayaan, sudut pandang, dan latar belakang, untuk memperkuat validitas dan generalisasi model.

\subsection{Preprocessing dan Augmentasi Data}
Sebelum digunakan untuk pelatihan, data gambar mengalami **proses preprocessing** untuk memastikan kualitas dan keseragaman gambar yang masuk ke dalam model. Langkah-langkah preprocessing meliputi:

\begin{itemize}
    \item **Resizing Gambar**: Setiap gambar diubah ukurannya menjadi **224x224 piksel** agar sesuai dengan input ukuran yang dibutuhkan oleh model, seperti ViT dan DeiT.
    \item **Normalisasi**: Gambar dinormalisasi menggunakan **mean** dan **standar deviasi** yang dihitung dari dataset **ImageNet** untuk memastikan keseragaman dalam distribusi warna dan intensitas pixel antar gambar. Nilai yang digunakan adalah: \(\text{mean} = [0.485, 0.456, 0.406]\) dan \(\text{std} = [0.229, 0.224, 0.225]\).
    \item **Konversi ke Tensor**: Gambar diubah menjadi format tensor agar dapat diproses oleh model berbasis PyTorch.
\end{itemize}

Untuk meningkatkan keragaman dan mengurangi overfitting, dataset juga mengalami **augmentasi data**, yang melibatkan berbagai transformasi acak pada gambar, seperti:

\begin{itemize}
    \item **Rotasi**: Gambar diputar secara acak dalam rentang tertentu untuk menambah variasi sudut pandang.
    \item **Flipping Horizontal**: Gambar dibalik secara horizontal untuk meningkatkan ketahanan model terhadap orientasi objek yang berbeda.
    \item **Zoom**: Zoom in/zoom out acak dilakukan untuk meningkatkan keberagaman pada variasi ukuran objek dalam citra.
\end{itemize}

Augmentasi ini diharapkan dapat memperbaiki kemampuan generalisasi model dengan menciptakan data tambahan yang lebih beragam.

\subsection{Konfigurasi Training}
Pelatihan model dilakukan dengan menggunakan **optimizer Adam** yang terkenal karena kemampuannya dalam mengoptimalkan kecepatan konvergensi pada masalah dengan banyak parameter. Untuk pengaturan awal, digunakan **learning rate sebesar 1e-4** yang dipilih berdasarkan eksperimen awal dan untuk memastikan pelatihan berjalan stabil. Dua model yang diuji dalam percobaan ini, yaitu **Vision Transformer (ViT)** dan **Data-efficient Transformer (DeiT)**, masing-masing dilatih selama **10 epoch** untuk memastikan bahwa model dapat belajar secara optimal dari data pelatihan.

Beberapa konfigurasi pelatihan yang diterapkan pada kedua model ini adalah sebagai berikut:
\begin{itemize}
    \item **Ukuran Batch**: Setiap iterasi pelatihan menggunakan batch berisi **32 gambar**. Ukuran batch ini dipilih untuk menjaga keseimbangan antara penggunaan memori dan stabilitas gradient.
    \item **Loss Function**: Fungsi kerugian yang digunakan adalah **CrossEntropyLoss**, yang umum digunakan untuk masalah klasifikasi multi-kelas. Fungsi ini mengukur perbedaan antara distribusi probabilitas yang diprediksi dan label target yang sebenarnya.
    \item **Epoch**: Kedua model, ViT dan DeiT, dilatih selama **10 epoch** untuk memastikan model memperoleh konvergensi yang memadai pada dataset pelatihan.
\end{itemize}

Proses pelatihan ini dilakukan pada **CPU** dengan spesifikasi perangkat keras yang terbatas, mengingat keterbatasan sumber daya yang ada. Untuk keperluan evaluasi, setiap epoch diukur berdasarkan **loss** dan **akurasi**, dengan tujuan untuk mengukur seberapa baik model dalam mengenali pola dan karakteristik citra makanan.

Selain itu, setelah pelatihan model, **checkpoint** disimpan pada setiap akhir epoch untuk memungkinkan pemulihan model jika pelatihan terganggu dan juga digunakan untuk evaluasi lebih lanjut setelah pelatihan selesai.

\subsection{Evaluasi Model}
Setelah pelatihan selesai, model diuji menggunakan dataset pengujian yang terdiri dari **270 gambar** yang belum pernah dilihat oleh model sebelumnya. Evaluasi dilakukan dengan menggunakan **loss function CrossEntropyLoss** dan mengukur **akurasi** dari prediksi model. Model juga dievaluasi berdasarkan **confusion matrix** untuk memeriksa distribusi kesalahan dalam klasifikasi, serta **visualisasi perbandingan** antara kedua model dalam hal **loss** dan **akurasi**. Tujuan evaluasi ini adalah untuk menentukan model mana yang lebih baik dalam hal **akurasi klasifikasi citra** dan **efisiensi komputasi**.

\subsection{Library dan Framework}
\begin{itemize}
    \item \textbf{Python}: 3.8+
    \item \textbf{PyTorch}: 2.0+ 
    \item \textbf{Torchvision}: 0.13+
    \item \textbf{Timm}: 0.4.12
    \item \textbf{Matplotlib}: 3.5+ (Untuk visualisasi)
    \item \textbf{Seaborn}: 0.11+ (Untuk visualisasi statistik)
    \item \textbf{Scikit-learn}: 1.0+ (Untuk evaluasi metrik seperti confusion matrix)
    \item \textbf{Pillow (PIL)}: 8.3+ (Untuk manipulasi gambar)
    \item \textbf{TensorBoard}: Untuk visualisasi pelatihan
    \item \textbf{Hugging Face Transformers}: Untuk model transformer lainnya
    \item \textbf{Optuna}: Untuk optimisasi hyperparameter
\end{itemize}


\subsection{Spesifikasi Hardware dan Software}
\begin{itemize}
    \item **Prosesor**: Intel Core i5
    \item **RAM**: 12 GB
    \item **GPU**: Tidak digunakan (hanya CPU)
    \item **Software**:
        \begin{itemize}
            \item **Python** 3.8+
            \item **PyTorch** 2.0+
            \item **Timm** 0.4.12
            \item **Torchvision** 0.13+
        \end{itemize}
\end{itemize}


\subsection{Cara Pengukuran Metrik Evaluasi}
Metrik evaluasi yang digunakan adalah:
\begin{itemize}
    \item \textbf{Akurasi:} Persentase gambar yang diklasifikasikan dengan benar.
    \item \textbf{Loss:} Nilai dari fungsi kerugian CrossEntropyLoss selama pelatihan.
\end{itemize}

\section{HASIL DAN ANALISIS}

\subsection{Perbandingan Jumlah Parameter}
\begin{itemize}
    \item \textbf{ViT:} Memiliki jumlah parameter lebih banyak dibandingkan dengan DeiT, yang mengarah pada penggunaan memori dan waktu pelatihan yang lebih besar. Model ViT memerlukan komputasi yang lebih tinggi dan waktu inferensi yang lebih lama.
    \item \textbf{DeiT:} Lebih efisien dalam hal jumlah parameter berkat penggunaan distillation, yang memungkinkan model bekerja lebih cepat dengan ukuran model yang lebih kecil, sehingga lebih hemat memori dan waktu pelatihan.
\end{itemize}

\subsection{Perbandingan Metrik Performa}
\begin{itemize}
    \item \textbf{Akurasi ViT:} 97.41\% (dengan hasil evaluasi pada dataset pengujian). ViT menunjukkan akurasi yang tinggi, namun memerlukan sumber daya yang lebih banyak.
    \item \textbf{Akurasi DeiT:} 87.11\% (pada pelatihan awal). Meskipun sedikit lebih rendah dibandingkan ViT, DeiT tetap menunjukkan kinerja yang baik meskipun menggunakan teknik distillation yang lebih efisien.
\end{itemize}

\subsection{Perbandingan Waktu Inferensi}
\begin{itemize}
    \item \textbf{ViT:} Lebih lambat dalam inferensi karena ukuran model yang lebih besar dan lebih banyak parameter, yang mempengaruhi kecepatan eksekusi.
    \item \textbf{DeiT:} Inferensi lebih cepat berkat optimisasi distillation, yang mengurangi jumlah parameter dan mempercepat proses pengolahan data.
\end{itemize}

\subsection{Perbandingan Plot Loss dan Akurasi}
Untuk lebih memahami perbandingan kinerja kedua model, berikut adalah **perbandingan plot** antara **loss** dan **akurasi** dari **ViT** dan **DeiT**. Plot ini menggambarkan perbandingan antara kedua model dalam hal **loss** dan **akurasi**, di mana **ViT** menunjukkan **akurasi yang tinggi** namun dengan **loss** yang sedikit lebih besar. Sementara itu, **DeiT** lebih efisien dalam hal **loss** dan tetap mencapai **akurasi** yang sangat baik meskipun dengan parameter yang lebih sedikit dan proses pelatihan yang lebih efisien.

\begin{figure}[h]
    \centering
    \includegraphics[width=0.7\textwidth]{comparison_plot.png}
    \caption{Perbandingan Plot Loss dan Akurasi antara ViT dan DeiT}
    \label{fig:comparison-plot}
\end{figure}

% Tempat untuk menambahkan gambar confusion matrix
\subsection{Confusion Matrix}
Confusion matrix memberikan gambaran lebih jelas tentang bagaimana masing-masing model (ViT dan DeiT) mengklasifikasikan gambar dalam lima kelas yang berbeda. Tabel ini menunjukkan **jumlah kesalahan** dalam pengklasifikasian setiap kelas, yang berguna untuk menilai **kekuatan model** dan **jenis kesalahan** yang mungkin terjadi.

\begin{figure}[h]
    \centering
    \includegraphics[width=0.7\textwidth]{confusion_matrix.png} % Path gambar yang benar
    \caption{Confusion Matrix untuk Model ViT dan DeiT}
    \label{fig:confusion-matrix}
\end{figure}

\subsection{Perbandingan Metrik Komprehensif}
Tabel berikut menunjukkan perbandingan metrik komprehensif antara ViT dan DeiT berdasarkan **akurasi**, **loss**, dan **waktu inferensi**. Meskipun ViT menunjukkan akurasi yang lebih tinggi, **DeiT** menunjukkan efisiensi yang lebih baik dalam hal waktu inferensi dan penggunaan sumber daya.

\begin{table}[h]
\caption{Perbandingan Metrik Komprehensif}
\label{tab:comprehensive-metrics}
\centering
\begin{tabular}{|l|l|l|}
\hline
\textbf{Metrik} & \textbf{ViT} & \textbf{DeiT} \\ \hline
Akurasi & 97.41\% & 87.11\% \\ \hline
Loss & 0.1466 & 0.3589 \\ \hline
Waktu Inferensi & Lebih lama & Lebih cepat \\ \hline
\end{tabular}
\end{table}


\subsection{Analisis Mendalam}
DeiT menunjukkan keunggulan dalam hal efisiensi komputasi, meskipun ViT memiliki akurasi yang lebih tinggi. Meskipun demikian, hasil perbandingan ini menunjukkan bahwa penggunaan DeiT sangat efektif dalam aplikasi dengan sumber daya terbatas.

\section{KESIMPULAN DAN SARAN}

\subsection{Kesimpulan Hasil Perbandingan}
Berdasarkan hasil eksperimen, ViT memiliki performa yang lebih tinggi dalam hal akurasi, tetapi membutuhkan sumber daya yang lebih besar. Sementara itu, DeiT menunjukkan efisiensi lebih baik dengan akurasi yang sedikit lebih rendah.

\subsection{Rekomendasi Model Berdasarkan Use Case}
Untuk aplikasi yang membutuhkan akurasi maksimal, ViT lebih disarankan, sementara untuk aplikasi dengan keterbatasan sumber daya, DeiT adalah pilihan yang lebih efisien.

\subsection{Saran untuk Pengembangan Lebih Lanjut}
Penelitian lebih lanjut dapat dilakukan untuk mengoptimalkan kedua model, dengan fokus pada peningkatan akurasi DeiT tanpa mengorbankan efisiensi komputasi.

\newpage
\section{LAMPIRAN}

\subsection{Informasi Repository GitHub}
\texttt{https://github.com/Joshia05/Tugas-deep-learning1.git} (masukkan link ke repository GitHub Anda).


\end{document}
